\section{Introduction}\label{sec:introduction}

The infinity Laplace equation is associated with minimizing the largest
local Lipschitz constant. For a smooth real function on an open subset of
$\R^d$, its unnormalized operator is
\[
 \Dinf u=\langle D^2u\,\nabla u,\nabla u\rangle
       =\sum_{i,j=1}^d u_i u_j u_{ij}.
\]
Only the second derivative in the gradient direction enters the equation.
This degeneracy makes control of the full gradient a central issue even
far from the boundary. We prove a uniform interior H\"older estimate for
the gradient, with a positive exponent depending only on the dimension.

\subsection{The result}
All balls are Euclidean; $B_r$ denotes the ball centered at the origin
unless a center is specified. A continuous function $u$ on an open set
$\Omega$ is \emph{infinity-harmonic} if every local upper $C^2$ test $\phi$
at $z\in\Omega$ satisfies $\Dinf\phi(z)\ge0$, and every local lower test
satisfies $\Dinf\phi(z)\le0$. All tests are included, also when
$\nabla\phi(z)=0$. Write $\osc_Eu=\sup_Eu-\inf_Eu$, and for a vector
field $g$ set
\[
 [g]_{C^{0,\alpha}(E)}
   =\sup_{\substack{x,y\in E\\x\ne y}}
       \frac{|g(x)-g(y)|}{|x-y|^\alpha}.
\]

\begin{theorem}[Uniform interior estimate]\label{thm:main}
For every integer $d\ge3$ there are constants $\alpha_d\in(0,1/3]$ and
$C_d<\infty$ such that every bounded infinity-harmonic function
$u\in C(B_1)$, where $B_1\subset\R^d$, belongs to
$C^{1,\alpha_d}(B_{1/2})$ and satisfies
\[
 \|\nabla u\|_{L^\infty(B_{1/2})}
       +[\nabla u]_{C^{0,\alpha_d}(B_{1/2})}
       \le C_d\osc_{B_1}u.
\]
\end{theorem}

The exponent arises from a compactness argument and is not explicit.
The qualitative interior regularity statement follows on any open domain.

\begin{corollary}[Interior $C^1$ regularity]\label{cor:interior-c1}
Let $d\ge3$, let $\Omega\subseteq\R^d$ be a nonempty connected open set,
and let $u\in C(\Omega)$ be infinity-harmonic. Then
$u\in C^1_{\mathrm{loc}}(\Omega)$.
\end{corollary}

\begin{proof}
For each $x_0\in\Omega$, choose $r>0$ with
$\overline B_{2r}(x_0)\subset\Omega$. Continuity bounds $u$ on this closed
ball, and $y\mapsto u(x_0+2ry)$ is infinity-harmonic on $B_1$.
Theorem~\ref{thm:main} applies and gives the conclusion near $x_0$.
\end{proof}

More precisely, the same change of variables gives
\[
 \|\nabla u\|_{L^\infty(B_r(x_0))}
   +r^{\alpha_d}[\nabla u]_{C^{0,\alpha_d}(B_r(x_0))}
   \le \frac{C_d}{r}\osc_{B_{2r}(x_0)}u,
\]
after enlarging $C_d$ if necessary.

\subsection{Context and methods}
Aronsson's work on absolute minimization connects the infinity Laplace
equation with a local version of the optimal Lipschitz extension problem
\cite{Aronsson1967}. A locally Lipschitz function $u$ is an absolute
minimizer if, for every open $V\Subset\Omega$ and every Lipschitz
competitor $w$ with $w=u$ on $\partial V$,
\[
  \|Du\|_{L^\infty(V)}\le\|Dw\|_{L^\infty(V)}.
\]
Aronsson identified the equation for smooth absolute minimizers.
Jensen established its viscosity characterization and uniqueness for the
Dirichlet problem \cite{Jensen1993}. Crandall, Evans, and Gariepy
developed the equivalent comparison-with-cones formulation and the
associated local slope estimates \cite{CrandallEvansGariepy2001}.
These results supply a robust theory of continuous solutions, while
regularity beyond local Lipschitz continuity presents a separate
difficulty: the equation controls second derivatives only in the gradient
direction.

The first-order regularity theory separates affine approximation from
uniqueness and continuity of the approximating slope. Building on the
cone and slope framework of Crandall, Evans, and Gariepy,
Crandall and Evans proved that infinitesimal blow-ups at a fixed point
are affine \cite{CrandallEvansGariepy2001,CrandallEvans2001}.
Savin proved continuity of the gradient in the plane
\cite{Savin2005}, and Evans and Savin obtained planar interior
$C^{1,\alpha}$ estimates for some positive universal exponent
\cite{EvansSavin2008}. Evans and Smart subsequently proved
differentiability at every point in every dimension
\cite{EvansSmart2011}. Their conclusion identifies a unique derivative
at each point; continuity of that derivative is a further assertion.
For a significant higher-dimensional special class, Peng, Zhang, and
Zhou proved interior $C^1$ regularity and nonvanishing of the gradient for
infinity-harmonic potentials in convex rings
\cite{PengZhangZhou2025}.

The planar endpoint has also been the subject of recent work. Brustad's
June 2026 preprint proves local $1/3$-H\"older continuity of the gradient
under the additional assumptions that it never vanishes and has a
Lipschitz inverse \cite[Theorem~1.1]{Brustad2026}.
Xu's July 2026 preprint states the unrestricted planar
$C^{1,1/3}_{\mathrm{loc}}$ theorem
\cite[Theorem~1.1]{Xu2026}; it separately leaves open a universal
endpoint norm estimate \cite[Remark~1.5(i)]{Xu2026}.
The theorem proved here concerns uniform interior estimates for arbitrary
solutions in dimensions $d\ge3$.

The restriction to an exponent at most $1/3$ is natural. Aronsson's
example \cite{Aronsson1984}
\[
  u(x)=|x_1|^{4/3}-|x_2|^{4/3},\qquad x\in\R^d,
\]
is infinity-harmonic and has
$|Du(he_1)-Du(0)|=\tfrac43|h|^{1/3}$.
It therefore rules out any exponent greater than $1/3$, in every
dimension $d\ge2$. Our compactness argument produces a positive exponent
depending on the dimension and does not identify the endpoint value.

\subsection{Proof strategy}
The proof has two parts. First we reduce failure of a uniform power
rate for affine approximation to an entire nonaffine solution of a
limiting equation. Then we prove that the estimates inherited by
that limit force it to be affine.

Normalize the original solutions to have oscillation at most one and
adjoin a linear coordinate, $U(y,b)=u(y)+b$. This preserves the equation
and ensures that accurate affine fits have nonzero slopes. A three-point
estimate controls the midpoint defect of a function close to an affine
function, divided by the endpoint separation, uniformly even at much
smaller separations. Combined with fixed-center affine blow-ups, it
gives affine approximation with error $o(r)$ uniformly across normalized
solutions and interior centers. This small-scale uniformity is essential:
compactness of the solutions alone would not give a common rate for
their derivatives.

We measure finer approximation on cylinders with axial radius $r$ and
transverse radius $\sqrt\lambda r$, allowing error $\lambda r$ from an
affine function whose slope points along the axis. A power rate of decay
for the worst fitting parameter makes consecutive slopes converge at a
uniform power rate, giving the gradient estimate. If every power bound
fails, a selection of almost-maximal normalized errors provides scales
with two properties: one fitting parameter fails at the selected center,
but a fixed multiple of it succeeds at every nearby center and every
fixed multiple of that scale. Those fits provide compactness after
stretching the transverse variables and magnifying the deviation from
the central affine fit.

The entire limit $v(x,t)$, with $x\in\R^d$ and $t\in\R$, satisfies
\[
 F[v]:=D^2v[(\nabla_xv,1)]=0,
 \qquad D^2v[\xi]=\xi^{\mathsf T}D^2v\,\xi.
\]
At every center $(x,t)$ and radius $R>0$ it admits coefficients
$p\in\R^d$ and $q\in\R$ such that
\[
 |v(x+h,t+\ell)-v(x,t)-p\cdot h-q\ell|\le C_0R
 \quad\text{if } |\ell|\le R,\quad |h-p\ell|\le R,
\]
with one finite constant $C_0$. The horizontal slope determines the
shear of the fitting box. Our Liouville Theorem states that every entire
solution with these fits is affine. An affine limit would improve the
selected failed fit, proving the desired power bound.

The anisotropic scaling and the limiting equation are due to Evans and
Savin \cite[Sections~2.3--2.4]{EvansSavin2008}. Their argument derives
gradient H\"older regularity from quantitative flatness estimates in thin
cylinders; the estimates required for that argument are established in
their paper in dimension two. Quadratic comparisons on horizontal slices
and monotone forward and backward slopes also appear in their analysis
\cite[Section~3]{EvansSavin2008}. Here uniform affine approximation,
comparison of intersecting cylinders, and a Liouville theorem for an
entire class with affine fits at every center and scale provide the
compactness argument. The Liouville proof combines an extremal
horizontal averaging defect, preservation of the subsolution inequality
by its variance-penalized envelope, and localization of contacts at
arbitrarily small times around a common terminal point.

For the Liouville argument, introduce, for $\ell>0$, the quadratic envelopes
\begin{align*}
 \Penv_\ell v(x,t)
   &=\sup_y\left\{v(y,t+\ell)-\frac{|y-x|^2}{2\ell}\right\},
 \\
 \Menv_\ell v(x,t)
   &=\inf_y\left\{v(y,t-\ell)+\frac{|y-x|^2}{2\ell}\right\}.
\end{align*}
The fits give subquadratic spatial growth, so the extrema are attained.
The quotients $(\Penv_\ell v-v)/\ell$ and
$(v-\Menv_\ell v)/\ell$ are nondecreasing in $\ell$ and have a common
small-length limit at each point. If both quotients were constant in
length at one point, quadratic comparison would force $v$ to be affine.
We seek such a limit while retaining an obstruction to affinity.

That obstruction is a horizontal Jensen defect: the excess of a weighted
average of values over the value at their spatial mean. Starting from one
nonaffine solution, we recenter and rescale configurations for which the
ratio of the defect to the root mean square distance from the mean
approaches its supremum.
Compactness then gives an extremal solution. For $t>0$, take the supremum
of these averages at time $t$ and a fixed spatial mean, after subtracting
their variance divided by $2t$. The extremal ratio bounds the resulting
barycentric envelope's excess over $v$ by a fixed multiple of time, with
equality at a selected point. Subtracting that linear function of time produces a subsolution
$W\le v$ touching $v$ at positive time; the unshifted envelope has a
strictly positive gap there. The variance penalty is affine along the
directions in the viscosity test, which makes the subsolution inequality
survive this averaging.

Forward optimizers for $W$ and backward optimizers for $v$ define two
possible steps from each space-time point. A fair choice between them
gives a drift inequality controlling the expected sum of squared discrepancies
of their velocities from a common reference velocity. This control is
unchanged by a fitted change of frame. It yields contacts at decreasing
positive times, with one terminal point remaining in a bounded part of
infinitely many fitted frames. The common small-length envelope slope at
that point is finite. After normalization at these shrinking contact
scales, every fixed envelope length corresponds to a length tending to zero
at the same terminal point. Differences of the normalized slopes cancel
the affine correction and tend to zero. A further limit therefore has
constant envelope slopes at every length, yet retains the positive
barycentric gap.
Rigidity gives the contradiction.

Section~\ref{sec:flatness} proves the lift, slope, and three-point estimates.
Section~\ref{sec:reduction} proves Theorem~\ref{thm:main} from the
Liouville Theorem. Sections~\ref{sec:limiting-tools}
and~\ref{sec:envelope-contact} construct the envelopes and extremal contact.
Sections~\ref{sec:discrete-propagation}
and~\ref{sec:contact-localization} propagate and localize that contact.
Section~\ref{sec:rigidity} proves rigidity and completes the Liouville
argument.

Throughout, equations with positive second-order sign use the same
viscosity convention as $\Dinf$: subsolutions satisfy the nonnegative
inequality at upper tests, and supersolutions the nonpositive inequality
at lower tests. Tests for equations in $(x,t)$ involve all variables
together. For a symmetric matrix $A$, we write $A[\xi]=\xi^{\mathsf T}A\xi$.
